\chapter{Podstawy teoretyczne}

Treść rozdziału pierwszego --- przegląd literatury i~podstawy teoretyczne dotyczące
tematu pracy.

\section{Podrozdział pierwszy}

Przykładowy tekst podrozdziału. Każda myśl zaczerpnięta od innego autora musi być
opatrzona przypisem~\cite{przyklad-ksiazka}.

Cytaty dosłowne należy ujmować w~cudzysłów: ,,Przykładowy cytat dosłowny z~publikacji
źródłowej''~\cite{przyklad-artykul}.

\section{Podrozdział drugi}

Nazwy łacińskie zapisujemy kursywą, np. \textit{Escherichia coli}.
Można również cytować materiały konferencyjne~\cite{przyklad-konferencja}
oraz źródła internetowe~\cite{przyklad-internet}.

Przykład odniesienia do tabeli --- patrz (Tab.~\ref{tab:przyklad}).

\begin{table}[htbp]
  \caption{Przykładowa tabela z~danymi}\label{tab:przyklad}
  \centering
  \begin{tabular}{lcc}
    \toprule
    \textbf{Parametr} & \textbf{Wartość A} & \textbf{Wartość B} \\
    \midrule
    Parametr 1 & 10 & 20 \\
    Parametr 2 & 30 & 40 \\
    Parametr 3 & 50 & 60 \\
    \bottomrule
  \end{tabular}
  \zrodlo{opracowanie własne.}
\end{table}
